% ─────────────────────────────────────────────────────────────────────────────
% RECOVERED SECTION — "Integrated interpretation of lipidome remodeling"
% Accidentally deleted in commit 2dfdac1 ("finaling", 2026-09-27).
% Last git copy was 1648ac4 (2026-09-26) but that version is STALE.
% This is the corrected text with all fixes applied. Insert into main.tex
% immediately after the Discussion opening paragraph and before
% \subsection*{Genetic control of lipid variation is stronger under LIN...}
% ─────────────────────────────────────────────────────────────────────────────

\subsection*{Integrated interpretation of lipidome remodeling across contrasting field trials}

The two most prominent changes in LIN are a decrease in SQDG and an increase in TG. The accumulation of TG is commonly associated with low temperatures, drought, oxidative stress, and nutrient-driven carbon reallocation, suggesting a shift towards increased storage within the LIN signature \citep{Tan2018, Yan2025, Pfaff2020}. The median LIN values are higher than CTL for every lipid ratio in which TG is the numerator, TG/SQDG (4.1$\times$), TG/MG (2.4$\times$), TG/MGDG (2.1$\times$), TG/DG (1.8$\times$), and TG/PE (1.7$\times$), so TG rises against every class it was tested against (Supplementary Table~S4, ratio statistics). DG rises less than TG but more than MG, giving DG/MG 1.3$\times$, while DG/PS falls to 0.32$\times$ because PS rises further still. Correlations also change for PS. PS--TG switches from negative in CTL ($r=-0.26$) to positive in LIN ($r=+0.43$), a strong sign reversal (Supplementary Table~S4, CLR correlation delta). Additionally, LINEX reactions add to these results. TG-producing branches shift towards products under LIN, while TG$\rightarrow$DG hydrolysis shifts towards substrate, collectively favoring TG over DG (Supplementary Table~S9, reaction balance). These class-level balances together suggest a reorganization of DG supply and TG synthesis and mobilization, though direct measures of enzyme activity are needed. Both arms contain LIN candidates. \textit{DGAT1} (\texttt{SORBI\_3010G170000}) acylates diacylglycerol to triacylglycerol and falls within a repeatedly mapped chromosome 10 crude-fat QTL in sorghum \citep{Boyles2017DGAT1, Yen2008DGAT}. The reverse (hydrolytic) step is \texttt{SORBI\_3001G041900}, an SDP1-clade lipase. In Arabidopsis, SDP1 hydrolyzes triacylglycerol to diacylglycerol and free fatty acid and accounts for most TG breakdown during reserve mobilization \citep{Kelly2011SDP1L}. This locus clusters with \textit{SbSDP1-L} in the sorghum lipase phylogeny \citep{Sapara2024PLA}.

Another change involves the balance among phospholipid headgroups. PS increased against every class it was tested against (Supplementary Table~S4, ratio statistics), and across genotypes the PS--PC correlation flipped from positive in CTL ($r=+0.45$) to negative in LIN ($r=-0.34$), the largest shift among class pairs, so the two headgroups move together across the panel under CTL but trade off against each other under LIN (Supplementary Table~S4, CLR correlation delta). PE showed the largest increase in mean double-bond number of any class, although the species detected differ slightly between trials, and this is a composition measure rather than a direct account of acyl editing. Comparable shifts in headgroup proportions, membrane unsaturation, and lysophospholipid pools are reported under cold, nutrient, oxidative, and water limitation, where they accompany changes in membrane fluidity, lipid repair, signaling, and remobilization in other studies \citep{Hou2016, Rawat2021, Yan2025, Zhang2020}. The balance between the lysophospholipid classes also shifted, with LPC/LPE dropping to 0.48-fold in LIN (Supplementary Table~S4, ratio statistics), driven by rising LPE while LPC remained essentially unchanged (Supplementary Table~S4, composition \%TIC). PS and the lysophospholipids each contribute under 0.2\% of TIC, so these are large proportional changes on small pools. The annotation is also shallow, with two PS species under CTL and one under LIN, four LPC species under CTL and three under LIN, and a single LPE species in both, so none of these shifts can be read as evidence about phospholipase activity.

These changes do not indicate phosphate conservation, nor a headgroup swap. Under phosphate limitation, phospholipids are degraded to release phosphate and are replaced by glycolipids, raising SQDG/DGDG and lowering phospholipids \citep{Essigmann1998, Hartel2000, Jouhet2004, KellyDormann2002, Pfaff2020}. However, neither occurs here. On the CLR scale, SQDG drops the most ($-0.942$, $0.39\times$), galactolipids drop less (MGDG $0.77\times$, DGDG $0.85\times$), and PG shows no change (Supplementary Table~S4), so SQDG falls while PG, the anionic thylakoid lipid it would normally replace, stays flat, and the two together fall from 3.13\% to 1.72\% of TIC. Phospholipids hold steady (30.3\% to 30.2\% of TIC) while glycolipids fall (47.5\% to 43.5\%), raising the ratio of phospholipids to glycolipids. Additionally, all phospholipids increase relative to SQDG (PC $2.07\times$, PE $2.41\times$, PG $2.54\times$; Supplementary Table~S4). PC salvage via deacylation would increase LPC, but LPC does not accumulate ($0.87\times$ CLR), so we find no evidence for that route behind the modest PC decline. Overall, the lipidome suggests a smaller plastid lipid fraction with a disproportionate loss of SQDG. Because SQDG is the only sulfur-containing membrane lipid, constrained sulfolipid synthesis is one explanation for why it falls further than the galactolipids, though we did not measure sulfur and other explanations remain open. The LIN GWAS returned one gene annotated as a sulfoquinovosyl transferase SQD2 (\texttt{SORBI\_3002G000600}) and one paralogue from the same glycosyltransferase family (\texttt{SORBI\_3003G076900}), both under fatty-acid traits rather than any SQDG trait, so they are candidates to follow up rather than evidence for the sulfolipid route. Without sulfur, chlorophyll, or thylakoid measurements the lipidome alone cannot separate these possibilities, and we therefore describe the LIN lipidome as a redistribution among plastid, phospholipid, lysophospholipid, and neutral-storage pools rather than a phosphate-driven remodeling.
